% TOP_PROBLEM: {"problemId": "problem.small-cohen-macaulay-conjecture", "canonicalTitle": "Small Cohen-Macaulay Conjecture", "releaseRank": 449, "primaryDomain": "algebra_representation_category", "primaryDomainLabel": "Algebra, representation & category theory", "displayStatus": "open_with_solved_subcases", "releaseStatus": "open_with_solved_subcases"}
% !TeX program = lualatex
% Definition and sources only; this file is not an LLM solution attempt.
\documentclass[11pt]{article}
\usepackage[margin=25mm]{geometry}
\usepackage{fontspec}
\setmainfont{Latin Modern Roman}
\usepackage{amsmath,amssymb,mathtools}
\usepackage{unicode-math}
\setmathfont{Latin Modern Math}
\usepackage{xurl,hyperref}
\hypersetup{colorlinks=true,urlcolor=blue}
\setcounter{secnumdepth}{0}
\setlength{\parindent}{0pt}
\setlength{\parskip}{5pt}
\setlength{\emergencystretch}{3em}
\begin{document}
\section*{\texorpdfstring{Small Cohen-Macaulay Conjecture}{Small Cohen-Macaulay Conjecture}}\label{449}
\subsection{Definitions and mathematical statement}
Let $(R,\mathfrak m)$ be a complete Noetherian local ring of Krull dimension $d$. A system of parameters is a $d$-tuple $x_1,\ldots,x_d$ generating an $\mathfrak m$-primary ideal. A nonzero finitely generated $R$-module $M$ is maximal Cohen--Macaulay if $\operatorname{depth}_R M=d$: equivalently a system of parameters is an $M$-regular sequence, meaning each $x_i$ acts injectively on $M/(x_1,\ldots,x_{i-1})M$ and $M/(x_1,\ldots,x_d)M\ne0$.

The conjecture asks whether every such $R$ admits a maximal Cohen--Macaulay module. The word small means finite generation; an infinitely generated big Cohen--Macaulay module does not answer this question. The requested object is a module, not necessarily a finite algebra or a faithful module.
\par\smallskip\textit{Formulation and concept references:} \href{https://sites.lsa.umich.edu/hochster/wp-content/uploads/sites/1337/2024/10/J-splitsfrcm.pdf}{[S1]}.

\subsection{Short English statement}
Does every complete Noetherian local ring have a nonzero finitely generated module of full possible depth?
\subsection{Sources}
\begin{itemize}
\item[S1]Hochster and Yao, Transactions of the AMS 376 (2023), Question 5.1 and the following reduction sentence.
\url{https://sites.lsa.umich.edu/hochster/wp-content/uploads/sites/1337/2024/10/J-splitsfrcm.pdf}.
\textit{Formulation source.}
\end{itemize}
\end{document}
